\section{The Past as Death and Life}

\begin{quotationx}
This is the second section of \textbf{Guido de Giorgio}`s essay The Instant and Eternity. Here de Giorgio compares modern and traditional man. He brings in many themes previously covered here: the continuity of the past; the notion of depth in the flow of history, i.e., ``subterranean history''; the past understood in the interiority, not as something external. Moderns can never rest, they can never ``be'', because they are always becoming, looking to the future to pass judgment on them. But the moderns of the future will despise them as much as the moderns of today despise their own roots, heritage and past.

Burying beetles bury the carcasses of small vertebrates such as birds and rodents as a food source for their larvae. This is the perfect description of modern man as a beetle. Although the essay was written in 1939, it is even more applicable today.
\end{quotationx}


Modern man can be compared to a burying beetle nécrophore that longs for the day that never comes: the cadaver that he carries is the past, his inert, sterile heritage, and the day that he awaits is the future, the imaginary descent, the glorious completion of a chimeric unfinished childbirth. We will remark that all modern men, the ``great men of history'', wait for a definitive judgment of the future on their oeuvre, for perhaps they feel, consciously or not, that nothing of what they made relates traditionally to the royal stream of the past, not is it capable of resisting the movements of the magnetic needle of the present, the brief and momentary instant having an impact on many abysses other than the insignificant trace of the passing cloud. That is why ancient man is a bearer of worlds: he did not leave the past behind him, but harvests it and carries it along, in such a way to construct in reality a single incidental point, the sole present, the current time, while modern man, discarding a heavy burden for his not very virile shoulders, is light, inconsistent, and through fear of being thrown to the earth by the blows of the crosswinds, clings to the machine, which is both his cradle and his tomb. For the myth of the future is associated with the myth of speed which---if we consider its function, its interior plan---is the abolition of the past in the already traversed, the imperceptibility of the present minimized in the permanent expectation of the future. The readers who would like to deepen these insights in a penetrating manner will find more than one way leading easily to the comprehension of some major truths. We desire to establish here, with a certain insistence, only some critical reflections whose development perspective will turn out to be clearer and surer.

We therefore understand that modern man and ancient man are absolutely opposed and like the antipodes, in the literal sense of the term, the one by relationship to the other: tied to a same line of descent but turned toward different heavens with different constellations, although the same impassible sun throws light on that line of descent in what for one is day for the other, night. For the Ancients, in effect, the past is everything, for the moderns, nothing, even when they have the illusion of absent mindedly looking in the past for solutions to the questions of the present---what are called the ``warnings'', the ``teachings'' of the past--, around sentimental fantasies exploited with a cynical opportunism according to the circumstances and proposed to the credulity of the naïve for the most pitiable perpetrations. Rhetoric, which triumphs today as never before in today's cloudy and swampy Europe, has recourse to the most bestial ruses to obtain the consent of the plebes as the audience and makes use of the past as a remedy against all illnesses, the universal balm, the supports of the present, but of a momentary usage as if to ward off the \emph{Vae soli}! Woe to him who is alone, Eccl. 4:10 Modern man, in reality, is already frozen in the past, no longer lives it, and takes from it only dust and ruin: he studies it, classifies it, ignores it. The more detailed the inquiry is made, the more it becomes sketchy, each one seeking subsequently to make life blow over these bones asleep in the slumber of death. Thus, when they turn toward the past to study it, the moderns then succumb to the same illusion as when they believe, for example, that the photograph is closer than the truth, although it is denatured totally in being fixing in the momentariness of something already passed. But independently of study, let us see if the moderns use the past according to life. Who says past, says tradition, i.e., interior, dynamic unification, not exterior adhesion, not opportunistic sympathy, not simple position or situation; in other words, there should be between the past and the present, \emph{continuity}, \emph{immutability}, or, better said, a rhythmic development so regular, permanent, internal, that it would appear to be indifferent. Antiquity, in fact, is characterized by a constant tonality which endures, immobile, from one epoch to the other; there is and must be a change, but it takes place in depth, in the interior strata, invisibly, we are tempted to say, in a way to not disrupt the regularity of the rhythm.


\flrightit{Posted on 2012-06-07 by Cologero}

